% Generated draft: requires mathematical and source review.
\documentclass[11pt]{article}
\usepackage[margin=25mm]{geometry}
\usepackage{fontspec}
\setmainfont{Latin Modern Roman}
\usepackage{amsmath,amssymb,mathtools,unicode-math}
\setmathfont{Latin Modern Math}
\usepackage{enumitem,needspace,xurl,hyperref}
\setlength{\parindent}{0pt}
\setlength{\parskip}{5pt}
\setlength{\emergencystretch}{3em}
\newcommand{\N}{\mathbb{N}}
\newcommand{\Z}{\mathbb{Z}}
\newcommand{\Q}{\mathbb{Q}}
\newcommand{\R}{\mathbb{R}}
\newcommand{\C}{\mathbb{C}}
\newcommand{\F}{\mathbb{F}}
\newcommand{\A}{\mathbb{A}}
\newcommand{\PP}{\mathbb{P}}
\newcommand{\E}{\mathbb{E}}
\newcommand{\eps}{\varepsilon}
\newcommand{\dd}{\,\mathrm{d}}
\begin{document}
\section*{UnsolvedMath ID: 10400169}
\textbf{Source problem code:} catalogue entry
\subsection{Definitions and mathematical statement}
Let \(\mathcal C\) be a semisimple ribbon category with finitely many isomorphism classes of simple objects, as in the source. Choose representatives \(X_i\). In the usual categorical-trace convention its \(S\)-matrix has entries given, up to the normalization convention, by
\[
 S_{ij}=\operatorname{tr}_{X_i\otimes X_j}
 \bigl(c_{X_j,X_i}\circ c_{X_i,X_j}\bigr),
\]
where \(c\) denotes the braiding. Let \(Z_{\mathrm{ss},\mathcal C}(M)\) denote the source's state-sum invariant, whenever its construction is defined, and let \(Z_{\mathrm{RT},\mathcal C}(M)\) denote the Reshetikhin--Turaev invariant in the nondegenerate setting.

\emph{Question (Y. Kawahigashi).} If \(S\) is not invertible, does \(Z_{\mathrm{ss},\mathcal C}\) still have a description analogous to the comparison available when \(S\) is invertible? The source describes the invertible case by
\[
 Z_{\mathrm{ss},\mathcal C}(M)
 =\left|Z_{\mathrm{RT},\mathcal C}(M)\right|^2.
\]
This is the source's motivating comparison, not a proposed identity involving an undefined ordinary Reshetikhin--Turaev invariant for singular \(S\). A satisfactory answer should specify an analogous invariant or categorical construction and its relationship to the state sum in the degenerate case.

The excerpt does not supply the coefficient field, finiteness assumptions on morphism spaces, global-dimension restrictions, normalization, or manifold and auxiliary-structure conventions needed for these constructions. In particular, the displayed absolute-value formula presupposes complex-valued invariants and appropriate reality or unitarity conventions that are not consequences of the explicitly stated hypotheses alone. The draft preserves that motivating formula with this qualification rather than asserting it for every semisimple ribbon category.
\subsection{Short English statement}
For a semisimple ribbon category with finitely many simple objects but singular \(S\)-matrix, Kawahigashi asks for a description of its state-sum invariant analogous to the usual nondegenerate comparison with Reshetikhin--Turaev theory. The exact replacement invariant and required hypotheses are part of the unresolved question.
\subsection{Sources}
\begin{itemize}
\item Local UnsolvedMath registry, permanent ID 10400169, source code catalogue entry; attribution Y. Kawahigashi. The supplied title is truncated after ``Problem 9.5 --- (Y.''.
\item Ohtsuki, \emph{Problems on invariants of knots and 3-manifolds} (2002), Problem 9.5, PDF page 139. Supplied URL: \url{https://msp.org/gtm/2002/04/gtm-2002-04-024s.pdf}; recorded access 2026-07-29; extraction pdf-font-structure. The collection distinguishes 231 numbered bold problem, conjecture, and question headings.
\item The supplied registry labels the item partially\_solved and PARTIAL-PROGRESS and says it is related to modular categories and the quantum double construction. It supplies no specific partial-result theorem or citation. Background status evidence remains NEEDS\_REVIEW pending a current primary-literature audit, and redistribution terms for the publicly accessible author/publisher copy require review.
\end{itemize}
\end{document}
